% Section content template for individual Educative lessons
% This file contains the actual content for a specific section
% Generated from Educative API section content

% Section: Key Insights and Tips: Designing an Online Blackjack Game
% Section ID: 5475450393329664
% Section Slug: key-insights-and-tips-designing-an-online-blackjack-game
% Generated: 2025-12-12T12:37:02.625333

You have completed the case study for designing an online Blackjack game. This lesson adds to your preparation for low-level design interview questions, especially those involving games like Blackjack. It will highlight common mistakes to avoid, provide important tips for interviews, offer a short quiz to test what you know, and suggest other case studies that can help you understand the ideas better.

\subsection{Common mistakes}\label{lzEVpKDWS5DGTyrwz_ylV}
\\
To build a better and bigger online Blackjack game, try to avoid these common mistakes:

\begin{itemize}
\item
\protect\phantomsection\label{_tLiyGUCfH-2KaCr9VKgS}
\textbf{Lack of clear game control flow:} Not establishing a central class to manage the overall game logic, from wagers to payouts. You should implement a \texttt{BlackjackGame} class to control the game flow, accept wagers, distribute cards, and update game status.
\item
\protect\phantomsection\label{AY-V3VvMc4plqA3twjzxq}
\textbf{Not addressing ties:} Failing to define what happens when the player and dealer have the same points. You should implement a rule allowing the player to take their bet back or replay the game in case of a tie.
\item
\protect\phantomsection\label{mPbRd1z3lCp9TOGazfrPS}
\textbf{Overlooking distinct player roles:} Not differentiating between the player and dealer roles with their specific actions and rules. You should define an abstract \texttt{Player} class with \texttt{BlackjackPlayer} and \texttt{Dealer} derived classes, each with distinct responsibilities like placing bets or dealing cards.
\item
\protect\phantomsection\label{j5_oRRwHKNiTD4HhtKhNC}
\textbf{Poor handling of game states:} Not clearly defining the various phases of the game, such as shuffling, dealing, hitting, and standing. You should apply the State design pattern to manage the game's finite states, ensuring correct transitions between them.
\item
\protect\phantomsection\label{2U5xfvo34NkHivhqcayvy}
\textbf{Ignoring player/dealer busting rules:} Allowing hands to continue beyond 21 points without ending the game. You must implement a clear rule: if a hand exceeds 21 points, the player or dealer busts and loses the game.
\end{itemize}

\subsection{Key interview tips}\label{N4gsxS0DzIIbVWtxcZXjh}
\\
The following table shows some tips to help you prepare for an interview where you're asked to design an online blackjack game:

\begin{center}
\begin{tabular}{|p{0.28\linewidth}|p{0.68\linewidth}|}
\hline
\textbf{Tip} & \textbf{Explanation} \\
\hline

Clarify requirements. &
Explicitly confirm the rules of Blackjack with your interviewer, especially card values, hitting rules for the player and the dealer, and winning conditions. This shows attention to detail and foundational understanding. \\
\hline

Identify core components of the game. &
Discuss logical classes like Card, Deck, Shoe, Hand, and Player early in your design. This shows a modular and object-oriented approach. \\
\hline

Consider extensibility for different card games. &
Briefly mention how your Card, Deck, and Shoe classes could be generic enough to support other card games with minor modifications. This demonstrates forward-thinking design. \\
\hline

Justify your choice of design patterns. &
Explain why patterns like Iterator (for drawing cards) and State (for game flow) suit this problem. This showcases understanding of clean, scalable design practices. \\
\hline

Be prepared to discuss scalability for multiple concurrent games. &
Even if the problem focuses on one player vs. dealer, be ready to explain how your design can scale to handle many players simultaneously—e.g., by running multiple game instances or using a session manager. \\
\hline

\end{tabular}
\end{center}

\subsection{Test your understanding}\label{2I7Mr5mfph8iFXCiPpDo3}
\\
Take this quick quiz to check your understanding of designing an online blackjack game:

\subsection*{Designing an Online Blackjack Game}
\vspace{0.3cm}
\noindent
\textbf{Question 1:} Which design pattern is suitable for assigning cards to players from the deck by iterating through a list of cards?
\\[0.2cm]
\begin{itemize}
 \item State design pattern
 \item Observer design pattern
 \item Strategy design pattern
 \item \textbf{[CORRECT]} Iterator design pattern
\\[0.1cm]
\textit{Explanation:} 
The Iterator design pattern can be applied since cards are assigned to players from the deck by iterating through the list of cards.
\end{itemize}
\vspace{0.3cm}
\noindent
\textbf{Question 2:} The ``Block member'' use case involves modifying a member's status. What relationship exists between ``Block member'' and ``Modify member'' use cases?
\\[0.2cm]
\begin{itemize}
 \item Include
 \item \textbf{[CORRECT]} Extend
\\[0.1cm]
\textit{Explanation:} 
The ``Block member'' use case has an extend relationship with the ``Modify member'' use case, as blocking membership specifically modifies the member status.
 \item Composition
 \item Association
\end{itemize}
\vspace{0.3cm}
\noindent
\textbf{Question 3:} What type of UML relationship exists between the \texttt{Player} and \texttt{Hand} classes?
\\[0.2cm]
\begin{itemize}
 \item Association
 \item Aggregation
 \item \textbf{[CORRECT]} Composition
\\[0.1cm]
\textit{Explanation:} 
The Player class is composed of Hand, indicating that a Hand is an integral part of a Player and typically doesn't exist independently.
 \item Inheritance
\end{itemize}

\subsection{Related case studies}\label{ZmhADjbi2ACsgH6q8t7iK}
\\
Explore these related case studies to deepen your understanding of the design principles applied in the online Blackjack game:

\begin{itemize}
\item
\textbf{Designing an online stock brokerage system:} This involves managing real-time transactions and account balances, dealing with high concurrency, similar to managing bets and payouts in a game.
\item
\textbf{Designing an online chess game:} This requires managing game states, player moves, and turns, and validating game rules, which parallels state management and rule enforcement in Blackjack.
\item
\textbf{Designing an online gaming platform (General):} This broader problem encompasses managing multiple concurrent games, player accounts, leaderboards, and real-time interactions, which are scalable extensions of the Blackjack game.
\item
\textbf{Designing a digital wallet/payment system:} This involves handling financial transactions, managing balances, and ensuring secure and accurate money transfers, directly relevant to betting and payouts.
\item
\textbf{Designing a fantasy sports league:} Requires managing player data, team compositions, scoring rules, and real-time updates based on external events, similar to tracking game state and player hands.
\end{itemize}

% End of section content